\section{Tema 2}
\subsection*{Transformaciones Estructurales}

\begin{enumerate}[label=\arabic*., series=preguntas]

    \item \textbf{El crecimiento de la renta per cápita en España generalmente no está acompañado de cambios en la estructura productiva.} \\
    \textbf{Falsa.} En el cuadro 1 se ve que la distribución de empleo entre actividades productivas ha variado drásticamente desde 1960-2023; la agricultura y pesca cayó del 38,7\% en 1960 al 3,6\% en 2023, los servicios aumentaron pasando del 31\% al 78,2\% y también hubo cambios en el caso de la industria, energía y construcción.

    \item \textbf{En las primeras fases de industrialización, la productividad del trabajo tiende a ser menor en la industria y los servicios que en la agricultura.} \\
    \textbf{Falsa.} De nuevo se ve lo mismo en el cuadro 1 (baja más en la agricultura, por matemáticas la productividad es mayor, aunque no se vea directamente). Además, en la pág. 18 de los apuntes se ve que hay un crecimiento medio anual de 4\% en cuanto al aumento de productividad (fase de industrialización).

    \item \textbf{A finales de los años noventa, el nivel de exposición de la economía española a la competencia exterior era significativamente inferior al de las economías europeas más avanzadas.} \\
    \textbf{Falsa.} No sale como tal, pero podemos verlo en el gráfico 9 donde se ve que el saldo exterior cae hacia lo negativo, lo que quiere decir que tenía una alta exposición y apertura comercial (déficit exterior por fuerte entrada de importaciones), desmintiendo que estuviera aislada de la alta competencia.

    \item \textbf{Una mayor competencia externa y la liberalización de los mercados interiores en España contribuyeron a un menor crecimiento económico.} \\
    \textbf{Falsa.} Hay un aumento del PIB real por encima del PIB potencial destacando los ciclos de 1986-1991, ya que se produce la entrada en la Comunidad Económica Europea y la apertura al exterior, y en 1995-2007 que se adopta el euro y mayor integración de mercados. En estos ciclos la economía de España se expandió de forma acelerada, no menor.

    \item \textbf{El aumento de las exportaciones generalmente disminuye el tamaño y la productividad de las empresas exportadoras.} \\
    \textbf{Falsa.} Desde el punto de vista lógico, ocurre justo lo contrario: se expanden para aprovechar las economías de escala y atender a mercados más grandes.

    \item \textbf{En 1960, la agricultura y la pesca representaban un porcentaje menor del empleo en España que en Alemania.} \\
    \textbf{Falsa.} Se puede ver en el cuadro 1 lo contrario: en 1960 la agricultura y la pesca concentraban el 38,7\% en España frente al 14\% en Alemania.

    \item \textbf{En 2021, el sector servicios empleaba a la mayoría de la población activa tanto en España como en Francia.} \\
    \textbf{Verdadera (con matiz en el año).} Aunque la tabla recoge datos para 2023 en lugar de 2021, la tendencia es clara. En 2023, el sector servicios suponía el 78,2\% del empleo en España y el 80,9\% en Francia.

    \item \textbf{La ampliación de los recursos públicos en España no ha tenido un impacto significativo en la productividad del sector privado.} \\
    \textbf{Falsa.} El documento indica que el sector público provee las infraestructuras y equipamientos que dan servicio al Capital Directamente Productivo (generado por el sector privado). Además, la inversión pública en educación genera ``Capital Humano'', el cual incrementa la productividad y la competitividad.

    \item \textbf{Las transferencias como las pensiones y el desempleo son una parte del gasto público que estimula el consumo.} \\
    \textbf{Verdadera/Justificación teórica:} Desde el punto de vista macroeconómico general, esto es correcto. Sin embargo, en estas diapositivas no hay ninguna mención explícita al impacto de las pensiones o el desempleo sobre el consumo, ya que las notas avisan de que los apuntes recogen solo cuadros y gráficos.

    \item \textbf{En las dos últimas décadas, España ha experimentado un avance significativo hacia una mayor equidad en la distribución de la renta y ha alcanzado una completa convergencia con las principales economías europeas en este aspecto.} \\
    \textbf{Falsa.} Se ha avanzado, pero no converge como tal; aún hay problemas de desigualdad en cuanto a criterios económicos. Además, estos términos no figuran en las diapositivas proporcionadas.

    \item \textbf{Una mayor equidad en la distribución de la renta puede favorecer el crecimiento económico al asegurar la vertebración social y la estabilidad institucional.} \\
    \textbf{Verdadera.} Aunque no hay referencias directas en las diapositivas a la vertebración social o su impacto en el crecimiento, conceptualmente es verdadera.

    \item \textbf{En 1990, la población inmigrante en España representaba un porcentaje significativo en comparación con otros países de la UE.} \\
    \textbf{Falsa.} En realidad solo era un 0,9\%. El texto explica que España tuvo un ascenso de la población en edad de trabajar ``algo más tardío'' y que se prolongó posteriormente debido a una ``posterior e intensa entrada de inmigrantes'', sugiriendo que la inmigración masiva llegó después de 1990.

    \item \textbf{La entrada de inmigrantes en España ha contribuido al aumento de la población en edad de trabajar y del empleo.} \\
    \textbf{Verdadera.} Se indica textualmente que el ascenso de la población en edad de trabajar en España se prolongó durante más tiempo ``debido al retraso en el baby boom y a una posterior e intensa entrada de inmigrantes''.

\end{enumerate}

\subsection*{Determinantes del Crecimiento a Largo Plazo}

\begin{enumerate}[label=\arabic*., resume=preguntas]

    \item \textbf{La tasa de variación de la renta per cápita se puede calcular restando las tasas de variación de la productividad del trabajo y del empleo per cápita.} \\
    \textbf{Falsa.} Se calcula mediante la \textit{suma} de las tasas de variación de ambas ratios, no restándolas.

    \item \textbf{El aumento de la proporción de población empleada no tiene límites.} \\
    \textbf{Falsa.} Matemáticamente, la fórmula del empleo per cápita ($\frac{Empleo}{Poblacion}$) tiene un límite físico. Está condicionada lógicamente a que no puede superar a la Población en Edad de Trabajar (PET) ni a la Población Activa (PA) disponibles en el país.

    \item \textbf{La productividad del trabajo es considerada clave para el crecimiento económico, ya que permite el incremento del salario real y de la renta familiar.} \\
    \textbf{Verdadera (conceptualmente).} El esquema vincula directamente el crecimiento económico con el ``Aumento de la renta per cápita'', y define la Productividad del trabajo como uno de los dos determinantes fundamentales.

    \item \textbf{Durante la década de 1990, la mayoría de los países europeos experimentaron una disminución en su capacidad de generación de empleo, similar a lo ocurrido en Estados Unidos.} \\
    \textbf{Falsa / No comprobable detalladamente.} El Cuadro 2 no desglosa la década de 1990. Sin embargo, en el periodo 1986-2024, el empleo per cápita creció en Estados Unidos (0,0\%) y en la mayoría de la UE creció en mayor o menor medida, no disminuyó.

    \item \textbf{Entre 1961 y 1985, la productividad del trabajo en España creció a una tasa menor que en Alemania.} \\
    \textbf{Falsa.} El Cuadro 2 muestra que en ese periodo la productividad del trabajo creció un 4,9\% en España, mientras que en Alemania lo hizo al 3,0\%. Por lo tanto, creció a una tasa mayor.

    \item \textbf{Entre 1986 y 2022, el empleo per cápita en España experimentó un crecimiento positivo.} \\
    \textbf{Verdadera.} El Cuadro 2 (periodo 1986-2024) muestra que el empleo per cápita en España tuvo un crecimiento positivo del 1,0\%.

    \item \textbf{La Población Activa (PA) es siempre mayor que la Población en Edad de Trabajar (PET).} \\
    \textbf{Falsa.} Por definición, la PA está compuesta por los miembros de la PET que están trabajando o buscando activamente empleo. Por tanto, la PA es un subgrupo y es siempre menor o igual que la PET.

    \item \textbf{El ascenso de la población en edad de trabajar fue más temprano e intenso en los países europeos menos desarrollados como España e Irlanda que en Estados Unidos.} \\
    \textbf{Falsa.} El texto afirma lo opuesto: ``El ascenso de la población en edad de trabajar fue más intenso en Estados Unidos que en la Europa Occidental durante los años 60, así como que, en los países europeos menos desarrollados, España e Irlanda en particular, ese ascenso fue algo más tardío...''.

    \item \textbf{En 2022, el porcentaje de población en edad de trabajar con respecto al total era mayor en España que en Alemania.} \\
    \textbf{Verdadera.} Según el Cuadro 3, en 2022 la población en edad de trabajar en España era el 65,4\%, frente al 63,2\% de Alemania.

    \item \textbf{Entre 1985 y 2022, el crecimiento del PIB per cápita en España se debió principalmente a un sustancial aumento de la productividad del trabajo.} \\
    \textbf{Falsa.} El Cuadro 2 muestra que entre 1986 y 2024 el PIBpc creció al 1,6\%; el aumento de la productividad fue solo del 0,6\%, mientras que el empleo per cápita aportó el 1,0\%. Hubo un estancamiento de los niveles de eficiencia y un nulo aumento de la PTF desde los años 90.

    \item \textbf{Desde 1970, la evolución del PIB per cápita y la productividad del trabajo en España han mostrado una convergencia notable.} \\
    \textbf{Falsa.} El Gráfico 3 muestra precisamente lo contrario. Ambas curvas se separan notablemente, con muchas oscilaciones inversas (si cae el PIB y se destruye empleo, la productividad sube, o viceversa).

    \item \textbf{En los periodos de recesión en España posteriores a 1986, el empleo por habitante ha aumentado, lo que ha provocado una disminución de la productividad del trabajo.} \\
    \textbf{Falsa.} Ocurre lo contrario. En las crisis (ej. 1993, 2009), la tasa de desempleo se dispara, destruyéndose empleo. Al destruirse empleo más rápido que la caída del PIB, la productividad (PIB/empleo) históricamente tiende a subir en recesiones (comportamiento anticíclico).

    \item \textbf{El comportamiento anticíclico de la productividad del trabajo durante la crisis de 2020 en España fue similar al de crisis anteriores.} \\
    \textbf{Falsa.} La crisis de 2020 fue una excepción. Las políticas (ERTE) sostuvieron el empleo, por lo que, al mantenerse el empleo mientras caía fuertemente la producción, la productividad se desplomó drásticamente en lugar de subir.

\end{enumerate}

\subsection*{Productividad, Capital y Progreso Tecnológico}

\begin{enumerate}[label=\arabic*., resume=preguntas]

    \item \textbf{El aumento de la productividad del trabajo se explica únicamente por el aumento del capital físico por trabajador.} \\
    \textbf{Falsa.} Según la función de Cobb-Douglas, la productividad del trabajo depende del capital físico por trabajador, del capital humano por trabajador y del nivel de eficiencia y progreso técnico (PTF).

    \item \textbf{La productividad total de los factores (PTF) mide la mejora en la eficiencia con la que se combinan trabajo y capital en el proceso productivo.} \\
    \textbf{Verdadera.} Se define como un multiplicador de la contribución al producto derivado de la combinación de factores, lo cual es función del progreso técnico.

    \item \textbf{La apertura al exterior y una mayor competencia en los mercados interiores tienden a disminuir la productividad total de los factores a corto plazo.} \\
    \textbf{Falsa.} Ocurre justo lo contrario.

    \item \textbf{El capital residencial se considera capital directamente productivo.} \\
    \textbf{Falsa.} El ``Capital Residencial'' es una categoría distinta al ``Capital Directamente Productivo'' (CDP) y se especifica que no se destina a la producción de bienes y servicios.

    \item \textbf{El capital humano incluye la salud, los conocimientos, las habilidades y los talentos de los trabajadores.} \\
    \textbf{Verdadera.} Se define exactamente como la combinación de salud, conocimientos, habilidades, destrezas y talentos útiles en la producción.

    \item \textbf{La inversión en el sistema educativo formal es una de las formas de generar capital humano.} \\
    \textbf{Verdadera.} Aparece explícitamente entre las formas de generar capital humano.

    \item \textbf{El gasto en educación en España como porcentaje del PIB ha experimentado un descenso constante desde la década de los 90 hasta la actualidad.} \\
    \textbf{Falsa.} Las conclusiones indican un aumento en el esfuerzo inversor, oscilando entre el 4\% y 5\% del PIB desde la década de los 90.

    \item \textbf{En 2022, el número medio de años de estudio de la población activa en España era similar al de 1960.} \\
    \textbf{Falsa.} En 2024 se situaba en 12,1 años, lo que supone más del doble que los 5,3 años registrados en 1960.

    \item \textbf{El aumento en el número de años de estudio en España se ha correspondido directamente con la evolución de la Productividad Aparente del Trabajo.} \\
    \textbf{Falsa.} El texto afirma literalmente que el aumento de años de estudio no se corresponde con la evolución de la Productividad Aparente del Trabajo.

    \item \textbf{La Productividad Aparente del Trabajo (PAT) en España se ha mantenido constante entre 1960 y 2022.} \\
    \textbf{Falsa.} Se registra un aumento total de +47.000€ en la PAT en el periodo 1960-2024.

    \item \textbf{En la función de producción Cobb-Douglas presentada, el parámetro $\alpha$ representa la elasticidad del producto respecto del capital humano.} \\
    \textbf{Falsa.} El parámetro $\alpha$ representa la elasticidad respecto del capital físico, mientras que $\beta$ representa la elasticidad respecto del capital humano.

    \item \textbf{La Productividad Total de los Factores (PTF) se calcula como la suma de la productividad media del trabajo, el capital físico por trabajador y el capital humano por trabajador.} \\
    \textbf{Falsa.} Se calcula como un residuo: restando a la variación de la productividad del trabajo las contribuciones del capital físico y humano ($\Delta A = \Delta y - \alpha\Delta k - \beta\Delta h$).

    \item \textbf{Entre 1960 y 2022, la mayor contribución al aumento de la PAT en España provino del aumento del capital humano.} \\
    \textbf{Falsa.} De los +47.000€ de aumento, el capital humano aportó solo el 15\%, mientras que la PTF aportó la mayor parte con un 49\%.

    \item \textbf{La PTF en España experimentó un notable aumento desde el inicio del decenio de 1990 hasta 2022.} \\
    \textbf{Falsa.} Se destaca un ``nulo aumento de la PTF desde el inicio del decenio de 1990'' debido al estancamiento de la eficiencia.

    \item \textbf{Desde el final del decenio de 1980 hasta la reciente crisis, la productividad del trabajo en España, en relación con la media de la UE-15, experimentó un aumento constante.} \\
    \textbf{Falsa.} La productividad del trabajo descendió desde el final de la década de 1980 hasta la Gran Recesión.

    \item \textbf{En las primeras etapas de desarrollo de un país, la importación de bienes de equipo no juega un papel relevante en el aumento del capital físico por trabajador.} \\
    \textbf{Falsa.} Sí juega, debido a que como estás en las primeras fases de desarrollo no desarrollas nada, lo más fácil es incorporarlos y supone un gran avance.

    \item \textbf{Un elevado grado de desarrollo facilita la adquisición y asimilación de tecnologías de otros países.} \\
    \textbf{Falsa.} El caso de España demuestra que no se ha sabido transitar.

    \item \textbf{En 2022, la ratio de intensidad investigadora en España superaba la de los países más desarrollados.} \\
    \textbf{Falsa.} Se indica que está algo por debajo.

    \item \textbf{El estancamiento de la PTF en España se puede atribuir principalmente a diferencias en el número de investigadores y en el nivel de educación de la población en comparación con otros países desarrollados.} \\
    \textbf{Falsa.} Dijimos que si consideramos estos dos indicadores como clave, en España han aumentado mucho, pero eso no se ha traducido en un aumento del progreso técnico.

    \item \textbf{España dedica un esfuerzo de inversión en I+D sobre el PIB similar al de Francia y Alemania.} \\
    \textbf{Falsa.} El esfuerzo inversor de España no es ``similar'' en absoluto, sino que está muy rezagado en comparación con el núcleo fuerte de Europa. Esta brecha en la inversión es una de las razones principales por las que a la economía española le cuesta tanto elevar su Productividad Total de los Factores (PTF) y asimilar o generar nuevo progreso tecnológico, tal y como se estudia al analizar las debilidades estructurales del país.

\end{enumerate}

\subsection*{Fluctuaciones Cíclicas y Desequilibrios Macroeconómicos}

\begin{enumerate}[label=\arabic*., resume=preguntas]

    \item \textbf{Los factores de oferta a largo plazo, como la población, el capital físico y humano y la PTF, no presentan distorsiones a corto plazo en cuanto a su evolución y ajuste de precios.} \\
    \textbf{Falsa.} Existen shocks transitorios de oferta y demanda que apartan a corto y medio plazo al PIB de su senda, creando el \textit{output gap}.

    \item \textbf{El PIB potencial describe la tendencia de la demanda agregada a largo plazo.} \\
    \textbf{Falsa.} Describe la tendencia de la \textit{oferta} agregada a largo plazo.

    \item \textbf{El output gap es la suma de las sendas de avance del PIB potencial y real.} \\
    \textbf{Falsa.} Es la \textit{diferencia} entre las sendas de avance del PIB potencial y real.

    \item \textbf{Durante el periodo 1995-2007, el PIB real en España se mantuvo por debajo de su PIB potencial.} \\
    \textbf{Falsa.} Ese periodo está catalogado expresamente bajo los ``aumentos del PIB real por encima del PIB potencial''.

    \item \textbf{Los principales momentos expansivos de la economía española han estado acompañados de disminuciones de precios.} \\
    \textbf{Falsa.} El Gráfico 8 muestra que la variación de precios se ha mantenido en valores positivos (inflación) a lo largo de las décadas de expansión, sin deflaciones marcadas.

    \item \textbf{La crisis del decenio de 1970 en España estuvo influenciada por shocks de oferta negativos, como las elevaciones del precio del crudo petrolífero.} \\
    \textbf{Verdadera.} Fue un shock de oferta (aunque en los gráficos solo se mencione genéricamente la Crisis de los años 70).

    \item \textbf{La masiva entrada de inmigrantes en los decenios de 1990 y 2000 constituyó un shock de oferta negativo para la economía española.} \\
    \textbf{Falsa.} La entrada de inmigrantes prolongó el ascenso de la población en edad de trabajar, lo cual es un factor determinante positivo para el crecimiento a largo plazo.

    \item \textbf{Tras el ingreso de España en las Comunidades Europeas y la entrada del euro, el rápido aumento de la demanda no generó tensiones alcistas sobre los precios.} \\
    \textbf{Falsa.} Trivial.

    \item \textbf{Antes de la entrada del euro, las devaluaciones de la moneda eran una herramienta utilizada para ganar competitividad y reestablecer el equilibrio exterior en España.} \\
    \textbf{Verdadera.} Cuando entró el euro perdimos esa posibilidad.

    \item \textbf{Entre 2002 y 2007, la reducción de los tipos de interés en la eurozona no tuvo un impacto significativo en la demanda interior española.} \\
    \textbf{Falsa.} Sí lo tuvo.

    \item \textbf{La crisis posterior a 2007 se caracterizó en España por un aumento de la demanda interna y una disminución de las exportaciones.} \\
    \textbf{Falsa.} Es justo lo contrario. El Gráfico 9 muestra que el ``Saldo exterior'' se disparó hacia terreno positivo tras esa crisis, lo que implica que las exportaciones superaron a las importaciones por el desplome de la demanda interna.

    \item \textbf{Desde la crisis de 2009, el sector exterior español ha mostrado un buen comportamiento a pesar de registrar incrementos del PIB superiores al 3\%.} \\
    \textbf{Verdadera.} El buen saldo se convierte en superávit.

    \item \textbf{El tardío proceso de apertura a la competencia internacional en España contribuyó a una menor intensidad de las fluctuaciones económicas en comparación con la media europea.} \\
    \textbf{Falsa.} Justo porque se tardó más se arrastran mayores fluctuaciones.

    \item \textbf{A partir de 2020, España experimentó un nuevo desequilibrio macroeconómico caracterizado por un incremento de la inflación.} \\
    \textbf{Verdadera.} El Gráfico 8 muestra un pronunciado y vertical repunte en la curva de ``Precios'' a partir de ese año.

\end{enumerate}
